\documentclass[10pt,a4paper]{amsart}
\usepackage[margin=19mm]{geometry}
\usepackage{amsmath,amssymb,mathtools}
\usepackage[scr=rsfs,cal=euler]{mathalfa}
\usepackage{xfrac}
\usepackage[unicode,hidelinks,bookmarks=false]{hyperref}
\newcommand{\ie}{i.e.\ }
\newcommand{\cf}{cf.\ }
\newcommand{\resp}{respectively\ }
\newcommand{\Subsection}{\subsection}
\providecommand{\nobreakdash}{\nobreak\hbox{-}}
\setlength{\emergencystretch}{2em}
\multlinegap0pt
\usepackage{amssymb}
\usepackage{xcolor}
\newtheorem{proposition}{Proposition}
\title{Persistence diagrams of random matrices via Morse theory: universality and a new spectral diagnostic}
\author{Matthew Loftus}
\date{Source: arXiv:2603.27903v1 (CC BY 4.0)}
\begin{document}
\maketitle

\noindent\emph{Source and licence.} Persistence diagrams of random matrices via Morse theory: universality and a new spectral diagnostic. Version arXiv:2603.27903v1 (CC BY 4.0), distributed by arXiv under the Creative Commons Attribution 4.0 International licence, http://creativecommons.org/licenses/by/4.0/, as recorded in the arXiv licence metadata for this exact version and shown on the abstract page https://arxiv.org/abs/2603.27903v1. Full licence notice: \texttt{licenses/arxiv-2603.27903-CC-BY-4.0.txt}.\par\smallskip

\noindent\textit{Editorial scope.} Counted unit: the article's proposition prop:pe\_formula (source0.tex lines 142--150). The article's complete original proof of it is delivered with it (source0.tex lines 152--160, one \texttt{proof} environment of 9 lines). No mathematics is written, repaired or re-derived: the statement and the proof are the article's own lines, in the article's own notation, and no retained line is substituted or re-flowed.  No further source-local context is needed: every cross-reference of every retained line resolves inside the two delivered spans.  Nothing was omitted from any retained span, so the package carries no omission record.  No retained line contains a citation, so the wrapper carries no bibliography and no bibliography entry.  No retained line contains a figure environment, an image-inclusion command or drawing code, so no image is delivered and the package declares no asset dependency.  Editorial formatting only, in addition: an AMS \texttt{amsart} preamble that restates the article's own declarations and macros used by the retained lines, namely the environments \texttt{proposition} on separate counters, together with the standard packages those lines need.  The retained environments print in this document as Proposition 1 in order of appearance, whereas the article prints the same items with the numbering of its own full document.  The complete native source of the article as distributed in the arXiv e-print for this exact version is bundled unchanged as \texttt{sources/197396/source0.tex}.
\begin{proposition}
\label{prop:pe_formula}
In the large-$n$ limit, the expected persistence entropy of an ensemble with density $\rho$ on $[\lambda_-, \lambda_+]$ is
\begin{equation}
\mathrm{PE} = \log(n \cdot \mathrm{TP}) + \frac{1}{\mathrm{TP}} \int_{\lambda_-}^{\lambda_+} \log \rho(\lambda)\, d\lambda + o(1),
\label{eq:pe_general}
\end{equation}
where $\mathrm{TP} = \lambda_+ - \lambda_-$.
\end{proposition}

\begin{proof}
Writing $p_k = s_k / \mathrm{TP} \approx 1/(n \cdot \mathrm{TP} \cdot \rho(\gamma_k))$, we have
\begin{align}
\mathrm{PE} &= -\sum_k p_k \log p_k \nonumber \\
&\approx -n \int \rho(\lambda) \cdot \frac{1}{n \cdot \mathrm{TP} \cdot \rho(\lambda)} \cdot \log \frac{1}{n \cdot \mathrm{TP} \cdot \rho(\lambda)}\, d\lambda \nonumber \\
&= \frac{1}{\mathrm{TP}} \int_{\lambda_-}^{\lambda_+} \left[\log(n \cdot \mathrm{TP}) + \log \rho(\lambda)\right] d\lambda \nonumber \\
&= \log(n \cdot \mathrm{TP}) + \frac{1}{\mathrm{TP}} \int_{\lambda_-}^{\lambda_+} \log \rho(\lambda)\, d\lambda.
\end{align}
\end{proof}

\end{document}
